% Solution: Tadelis (2013) Exercise 13.2 "Complete Information". Input by \solitem. Check: check_E1.py.
\begin{exsol}{Tadelis Exercise 13.2}{285}{Foundations}{2}
\asked $N$ bidders in a second-price auction with complete information: (a) is bidding one's value
still weakly dominant? (b) are there other Nash equilibria?

\method (a) The dominance argument of item~5 never used the others' information. (b) Look
for profiles where the winner is not the highest-value bidder, using that the winner pays the
second-highest bid (Tadelis \S13.1.1).

\solution
\textbf{(a) Yes.} The state-by-state comparison with the highest rival bid $r$ holds for every
$r$ whatever the bidders know, so $b_i=v_i$ is weakly dominant (items 5--6); in particular
$(v_1,\dots,v_N)$ is a Nash equilibrium. With $v_1>v_2\ge\dots$, bidder~1 wins and pays $v_2$.

\textbf{(b) Yes, many.} Take any bidder $k$ and let
$b_k=\max_{j\ne k}v_j$, $b_j=0$ for $j\ne k$ ($N\ge2$). Bidder $k$ wins and pays $0$, earning $v_k\ge0$;
he cannot do better and has no reason to deviate. Any other bidder $j$ can win only by bidding
$>b_k$ and then pays $b_k=\max_{l\ne k}v_l\ge v_j$, earning $v_j-b_k\le0$. So no one deviates:
a Nash equilibrium in which bidder~$k$ wins even if his value is lowest. Example, $v=(3,2,1)$: the
profile $(0,3,0)$ is an equilibrium with the value-$2$ bidder winning at price $0$; so is $(0,0,3)$
with the value-$1$ bidder winning.

\interp These equilibria use weakly dominated strategies ($b_k>v_k$ is weakly dominated by $b_k=v_k$)
and rest on the other bidders' staying out; under complete information nothing in Nash
equilibrium rules them out. Eliminating weakly dominated strategies selects the truthful
equilibrium, as in item~5.

\verify \texttt{check\_E1.py} enumerates all profiles on the integer bid grid $\{0,\dots,4\}^3$
for $v=(3,2,1)$ (fair coin at ties): the truthful profile and $(0,3,0)$, $(0,0,3)$ are Nash; 55
pure equilibria in all, in which every bidder wins in some equilibrium; dominance of
$b_i=v_i$ is verified against every pair of rival bids. Tadelis gives no answer.

\answer{(a) Yes, bidding the value is still weakly dominant. (b) Yes: e.g.\ one bidder bids the
highest rival value and everyone else bids $0$; that bidder wins at price $0$, whoever he is.}
\end{exsol}
